\section{习题}

\subsection{本章速览}

\begin{itemize}[leftmargin=1.6em]
  \item \textbf{定义}：$n$ 阶行列式是“取自不同行、不同列的 $n$ 个元素之积”的带符号和，共 $n!$ 项，符号由列下标排列的逆序数决定。对角线法则只适用于二、三阶，四阶以上必须用性质或展开。
  \item \textbf{六条性质}：转置不变 $|A^{T}|=|A|$；交换两行（列）变号；某行（列）乘 $k$ 则整个行列式乘 $k$；两行（列）相同或成比例则行列式为 $0$；对某一行（列）有可加性；把一行的 $k$ 倍加到另一行，行列式不变。注意行列式只是“对每一行线性”，一般 $|A+B|\neq|A|+|B|$。
  \item \textbf{展开定理}：$|A|=\sum_{j=1}^{n}a_{ij}A_{ij}$，其中代数余子式 $A_{ij}=(-1)^{i+j}M_{ij}$ 只与划去第 $i$ 行、第 $j$ 列后剩下的元素有关。用别的行的代数余子式去加权本行元素，和为 $0$。实际计算时先用倍加变换造零，再挑零最多的行（列）展开。
  \item \textbf{范德蒙德行列式}：$V_n=\prod_{1\le i<j\le n}(x_j-x_i)$。它为零当且仅当 $x_1,\dots,x_n$ 中有两个相等；识别题型的关键是各行（列）依次是同一个数的 $0,1,2,\dots$ 次幂。
  \item \textbf{分块与乘法}：分块三角 $\begin{vmatrix}A&C\\0&B\end{vmatrix}=|A||B|$；$|AB|=|A||B|$，从而 $|kA|=k^{n}|A|$，$|A^{*}|=|A|^{n-1}$，$|A^{-1}|=|A|^{-1}$。两块都在副对角线上的情形要额外记录换块产生的符号 $(-1)^{mn}$。
  \item \textbf{克莱姆法则}：$n$ 个方程、$n$ 个未知数，若系数行列式 $D\neq0$，则方程组有唯一解 $x_j=D_j/D$，$D_j$ 是把第 $j$ 列换成常数列所得的行列式。特别地，齐次方程组有非零解 $\iff D=0$。
\end{itemize}

\subsection{基础过关题}

\textbf{1.} 设 $A,B$ 均为 $n$ 阶方阵（$n\ge 2$），$k$ 为常数。判断下列命题是否正确，并说明理由：\\
（1） $|A+B|=|A|+|B|$；\quad （2） $|AB|=|A|\,|B|$；\quad （3） $|kA|=k|A|$；\quad （4） $|A^{T}|=|A|$；\quad （5） 若 $|A|=0$，则 $A$ 必有两行成比例。

\ifshowsolutions%
\boxed{\textbf{解}}
\textbf{思路：} 逐条对照六条基本性质。凡是“看起来像线性”的等式，都要警惕行列式只对单独一行线性；凡是否定性质的命题，用二阶反例即可推翻。

（1） \textbf{错}。取 $n=2$，$A=I_2$，$B=-I_2$。左边 $|A+B|=|O|=0$，右边 $|A|+|B|=1+1=2$。行列式的可加性只针对某一行（列），不是对整个矩阵加法成立的。

（2） \textbf{对}。这是行列式的乘法公式 $|AB|=|A||B|$，对任意同阶方阵成立。

（3） \textbf{错}。$kA$ 是把 $A$ 的每一行都乘以 $k$，提公因子要做 $n$ 次，故 $|kA|=k^{n}|A|$。取 $A=I_2$，$k=2$：$|2I_2|=4\neq 2|I_2|=2$。

（4） \textbf{对}。转置不改变行列式，这是最基本的性质。

（5） \textbf{错}。行列式为零只说明三行（列）线性相关，并不要求某两行成比例。例如
\[
A=\begin{pmatrix}1&2&3\\4&5&6\\7&8&9\end{pmatrix},
\]
第三行等于 $2r_2-r_1$（即 $r_3=2r_2-r_1$），故 $|A|=0$，但任意两行都不成比例。

\boxed{\textbf{答案：}\text{（2）（4） 正确，（1）（3）（5） 错误}}

\else \answerarea[6] \fi%

\textbf{2.} 计算四阶行列式
\[
D=\begin{vmatrix}
1&2&-1&3\\
2&5&1&0\\
3&1&2&4\\
-1&0&3&2
\end{vmatrix}.
\]

\ifshowsolutions%
\boxed{\textbf{解}}
\textbf{思路：} 用倍加变换把第一列下方元素化为零，再依次处理第二、三列，化成上三角。倍加变换不改变行列式的值，最后把主对角线元素相乘。

\begin{align*}
D&=\begin{vmatrix}
1&2&-1&3\\
2&5&1&0\\
3&1&2&4\\
-1&0&3&2
\end{vmatrix}
\xrightarrow{r_2-2r_1,\ r_3-3r_1,\ r_4+r_1}
\begin{vmatrix}
1&2&-1&3\\
0&1&3&-6\\
0&-5&5&-5\\
0&2&2&5
\end{vmatrix}\\
&\xrightarrow{r_3+5r_2,\ r_4-2r_2}
\begin{vmatrix}
1&2&-1&3\\
0&1&3&-6\\
0&0&20&-35\\
0&0&-4&17
\end{vmatrix}
\xrightarrow{r_4+\frac{1}{5}r_3}
\begin{vmatrix}
1&2&-1&3\\
0&1&3&-6\\
0&0&20&-35\\
0&0&0&10
\end{vmatrix}\\
&=1\times1\times20\times10=200.
\end{align*}
第三步中 $-4+\frac{1}{5}\times20=0$，$17+\frac{1}{5}\times(-35)=10$，故确实化为上三角。

\boxed{\textbf{答案：}D=200}

\else \answerarea[8] \fi%

\textbf{3.} 设
\[
D=\begin{vmatrix}
1&2&0&1\\
0&1&3&1\\
2&-1&3&0\\
2&1&0&3
\end{vmatrix},
\]
$A_{ij}$ 是元素 $a_{ij}$ 的代数余子式。求
\[
S_1=A_{31}+A_{32}+A_{33}+A_{34},\qquad
S_2=A_{31}-A_{32}+A_{33}-A_{34}.
\]

\ifshowsolutions%
\boxed{\textbf{解}}
\textbf{思路：} 代数余子式 $A_{3j}$ 是划掉第 $3$ 行、第 $j$ 列后算出来的，与第 $3$ 行本身的元素无关。因此把第 $3$ 行整体替换为 $(1,1,1,1)$ 后再按第 $3$ 行展开，得到的恰好是 $S_1$；替换为 $(1,-1,1,-1)$ 则得到 $S_2$。这类题不必逐个求余子式。

计算 $S_1$：把第 $3$ 行换成 $(1,1,1,1)$，再用第 $3$ 行把第 $1$ 列的其他元素消成 $0$，
\[
S_1=\begin{vmatrix}
1&2&0&1\\
0&1&3&1\\
1&1&1&1\\
2&1&0&3
\end{vmatrix}
\xrightarrow{r_1-r_3,\ r_4-2r_3}
\begin{vmatrix}
0&1&-1&0\\
0&1&3&1\\
1&1&1&1\\
0&-1&-2&1
\end{vmatrix}.
\]
按第 $1$ 列展开，只有 $a_{31}=1$，符号为 $(-1)^{3+1}=+1$：
\[
S_1=\begin{vmatrix}
1&-1&0\\
1&3&1\\
-1&-2&1
\end{vmatrix}
=1\cdot(3+2)-(-1)\cdot(1+1)+0=5+2=7.
\]
这里用的是三阶行列式按第一行展开：$1\cdot\bigl(3\cdot1-1\cdot(-2)\bigr)-(-1)\cdot\bigl(1\cdot1-1\cdot(-1)\bigr)$。

计算 $S_2$：把第 $3$ 行换成 $(1,-1,1,-1)$，
\[
S_2=\begin{vmatrix}
1&2&0&1\\
0&1&3&1\\
1&-1&1&-1\\
2&1&0&3
\end{vmatrix}
\xrightarrow{r_1-r_3,\ r_4-2r_3}
\begin{vmatrix}
0&3&-1&2\\
0&1&3&1\\
1&-1&1&-1\\
0&3&-2&5
\end{vmatrix},
\]
仍按第 $1$ 列展开：
\[
S_2=\begin{vmatrix}
3&-1&2\\
1&3&1\\
3&-2&5
\end{vmatrix}
=3(15+2)-(-1)(5-3)+2(-2-9)=51+2-22=31.
\]

\textbf{评注：} 原行列式 $|D|=45$，它与 $S_1,S_2$ 无关。若题目问的是 $a_{31}A_{31}+a_{32}A_{32}+a_{33}A_{33}+a_{34}A_{34}$，那才等于 $|D|$。

\boxed{\textbf{答案：}S_1=7,\ S_2=31}

\else \answerarea[8] \fi%

\textbf{4.} 计算
\[
V=\begin{vmatrix}
1&1&1&1\\
-2&0&1&3\\
4&0&1&9\\
-8&0&1&27
\end{vmatrix}.
\]

\ifshowsolutions%
\boxed{\textbf{解}}
\textbf{思路：} 逐行观察：第 $2$ 行是各列的 $1$ 次幂，第 $3$ 行是 $2$ 次幂，第 $4$ 行是 $3$ 次幂，第 $1$ 行是 $0$ 次幂，符合范德蒙德行列式的形状，其中 $x_1=-2,\ x_2=0,\ x_3=1,\ x_4=3$。

由范德蒙德公式
\[
V=\prod_{1\le i<j\le 4}(x_j-x_i),
\]
把六对差值逐一写出：
\[
\begin{aligned}
x_2-x_1&=0-(-2)=2, & x_3-x_1&=1-(-2)=3, & x_4-x_1&=3-(-2)=5,\\
x_3-x_2&=1-0=1, & x_4-x_2&=3-0=3, & x_4-x_3&=3-1=2.
\end{aligned}
\]
故
\[
V=2\times3\times5\times1\times3\times2=180.
\]

\textbf{评注：} 第 $2$ 列只有 $a_{12}=1$ 非零，若按第 $2$ 列展开，只剩一项
\[
V=(-1)^{1+2}\begin{vmatrix}-2&1&3\\4&1&9\\-8&1&27\end{vmatrix},
\]
右端三阶行列式中，每一行是 $(x,x^{2},x^{3})$，相当于范德蒙德行列式各列都乘了对应的 $x$，因此其值为
\[
x_1x_2x_3\prod_{i<j}(x_j-x_i)=(-2)\times1\times3\times30=-180,
\]
故 $V=180$，与上面一致。直接认出范德蒙德结构可以少走弯路。

\boxed{\textbf{答案：}V=180}

\else \answerarea[6] \fi%

\textbf{5.} 用克莱姆法则解线性方程组
\[
\begin{cases}
2x_1+x_2=1,\\
x_1+2x_2+x_3=1,\\
x_2+2x_3=3.
\end{cases}
\]

\ifshowsolutions%
\boxed{\textbf{解}}
\textbf{思路：} 先算系数行列式 $D$，确认 $D\neq0$（否则克莱姆法则不适用）；再把常数列分别替换第 $1,2,3$ 列，得 $D_1,D_2,D_3$，最后 $x_j=D_j/D$。

系数行列式（按第一行展开）：
\[
D=\begin{vmatrix}
2&1&0\\
1&2&1\\
0&1&2
\end{vmatrix}
=2\begin{vmatrix}2&1\\1&2\end{vmatrix}-1\begin{vmatrix}1&1\\0&2\end{vmatrix}+0
=2\times3-1\times2=4\neq0.
\]
所以方程组有唯一解。逐个替换列：
\[
\begin{aligned}
D_1&=\begin{vmatrix}1&1&0\\1&2&1\\3&1&2\end{vmatrix}
=1(4-1)-1(2-3)+0=3+1=4,\\
D_2&=\begin{vmatrix}2&1&0\\1&1&1\\0&3&2\end{vmatrix}
=2(2-3)-1(2-0)+0=-2-2=-4,\\
D_3&=\begin{vmatrix}2&1&1\\1&2&1\\0&1&3\end{vmatrix}
=2(6-1)-1(3-0)+1(1-0)=10-3+1=8.
\end{aligned}
\]
于是
\[
x_1=\frac{D_1}{D}=1,\qquad x_2=\frac{D_2}{D}=-1,\qquad x_3=\frac{D_3}{D}=2.
\]
代回验证：$2\times1+(-1)=1$；$1+2\times(-1)+2=1$；$(-1)+2\times2=3$，全部成立。

\boxed{\textbf{答案：}x_1=1,\ x_2=-1,\ x_3=2}

\else \answerarea[8] \fi%

\subsection{综合提高题}

\textbf{6.} 设 $A$ 是 $m$ 阶方阵，$B$ 是 $n$ 阶方阵。\\[2pt]
（1）证明 $\begin{vmatrix}0&A\\B&0\end{vmatrix}=(-1)^{mn}|A|\,|B|$；\\[2pt]
（2）利用（1）计算 $M=\begin{pmatrix}0&0&1&2\\0&0&3&4\\5&6&0&0\\7&8&0&0\end{pmatrix}$ 的行列式，并直接用初等变换复核。

\ifshowsolutions%
\boxed{\textbf{解}}
\textbf{思路：} 分块行列式的“零块在副对角线”情形不能直接套用分块三角公式，要先用列交换把 $A,B$ 换到主对角线上，而每次交换两列都会使行列式变号。

（1）把前 $n$ 列整体移到后 $m$ 列的右侧。这一步可以看成：第 $1$ 列依次与后面 $m$ 列相邻对换，共 $m$ 次；第 $2$ 列同样跨过 $m$ 列，共 $m$ 次……前 $n$ 列每列都跨过 $m$ 列，共 $mn$ 次相邻对换，每次变号，故
\[
\begin{vmatrix}0&A\\B&0\end{vmatrix}=(-1)^{mn}\begin{vmatrix}A&0\\0&B\end{vmatrix}=(-1)^{mn}|A|\,|B|.
\]
最后一步用了分块对角行列式等于各块行列式之积（先按前 $m$ 行展开，或反复用分块三角公式）。

（2）这里 $A=\begin{pmatrix}1&2\\3&4\end{pmatrix}$，$B=\begin{pmatrix}5&6\\7&8\end{pmatrix}$，$m=n=2$。$|A|=4-6=-2$，$|B|=40-42=-2$，$mn=4$，所以
\[
|M|=(-1)^{4}\times(-2)\times(-2)=4.
\]
直接复核：交换第 $1$ 列与第 $3$ 列、第 $2$ 列与第 $4$ 列（共 $2$ 次换列，符号不变），得
\[
|M|=\begin{vmatrix}1&2&0&0\\3&4&0&0\\0&0&5&6\\0&0&7&8\end{vmatrix}
=|A|\,|B|=(-2)(-2)=4,
\]
与公式一致。这里 $m=n=2$，恰好 $mn$ 为偶数，所以两种算法都得到 $+4$。

\boxed{\textbf{答案：}|M|=4}

\else \answerarea[8] \fi%

\textbf{7.} 计算 $n$ 阶行列式
\[
D_n=\begin{vmatrix}
a&1&1&\cdots&1\\
1&a&1&\cdots&1\\
1&1&a&\cdots&1\\
\vdots&\vdots&\vdots&\ddots&\vdots\\
1&1&1&\cdots&a
\end{vmatrix}\quad (a\text{ 为常数}).
\]

\ifshowsolutions%
\boxed{\textbf{解}}
\textbf{思路：} 除对角线外全是 $1$，把第 $n$ 行作为“基准行”，用其余各行减去它，就能把大片 $1$ 消成 $0$；再把前 $n-1$ 列加到最后一列，把剩余的非对角块消掉。

记第 $i$ 行为 $r_i$。对 $i=1,2,\dots,n-1$ 作 $r_i-r_n$：第 $i$ 行变为“第 $i$ 列为 $a-1$，其余前 $n-1$ 列为 $0$，最后一列为 $1-a$”，而 $r_n=(1,1,\dots,1,a)$。于是
\[
D_n=\begin{vmatrix}
a-1&0&\cdots&0&1-a\\
0&a-1&\cdots&0&1-a\\
\vdots&\vdots&\ddots&\vdots&\vdots\\
0&0&\cdots&a-1&1-a\\
1&1&\cdots&1&a
\end{vmatrix}.
\]
再把第 $1,2,\dots,n-1$ 列都加到第 $n$ 列上。对第 $i$ 行（$i<n$），前 $n-1$ 个元素中只有第 $i$ 列是 $a-1$，其余为 $0$，所以第 $n$ 列的元素变为
\[
(1-a)+(a-1)=0 ;
\]
对第 $n$ 行，前 $n-1$ 个元素都是 $1$，所以第 $n$ 列的元素变为
\[
a+\underbrace{1+\cdots+1}_{n-1\text{ 个}}=a+n-1 .
\]
于是行列式化为分块三角形状
\[
D_n=\begin{vmatrix}
a-1&0&\cdots&0&0\\
0&a-1&\cdots&0&0\\
\vdots&\vdots&\ddots&\vdots&\vdots\\
0&0&\cdots&a-1&0\\
1&1&\cdots&1&a+n-1
\end{vmatrix}
=(a-1)^{n-1}(a+n-1).
\]
最后一步：左上块是 $(a-1)I_{n-1}$，用分块三角公式，或按前 $n-1$ 行依次展开。

\textbf{评注：} 当 $a=1$ 且 $n\ge2$ 时，矩阵全为 $1$，公式给出 $(a-1)^{n-1}(a+n-1)=0$，符合秩为 $1$ 的事实；$n=1$ 时公式给出 $a$，也与 $D_1=a$ 一致。可用 $n=3,a=2$ 检验：$D_3=1^2\times4=4$，直接计算 $\begin{vmatrix}2&1&1\\1&2&1\\1&1&2\end{vmatrix}=4$。

\boxed{\textbf{答案：}D_n=(a-1)^{n-1}(a+n-1)}

\else \answerarea[10] \fi%

\textbf{8.} 设 $A$ 是 $3$ 阶方阵，$|A|=2$，$A^{*}$ 为 $A$ 的伴随矩阵。求\\
（1）$|A^{*}-3A^{-1}|$；\quad（2）$|(A^{*})^{-1}+A|$；\quad（3）$\begin{vmatrix}0&A\\A^{*}&0\end{vmatrix}$（这是一个 $6$ 阶行列式）。

\ifshowsolutions%
\boxed{\textbf{解}}
\textbf{思路：} 所有结论都来自两条基本关系：$AA^{*}=|A|I$（从而 $A^{-1}=A^{*}/|A|$、$|A^{*}|=|A|^{n-1}$）与乘法公式 $|AB|=|A||B|$。先把所求矩阵化成同一个矩阵的常数倍，再取行列式。

由 $|A|=2$，$n=3$，得
\[
A^{-1}=\frac{A^{*}}{|A|}=\frac{A^{*}}{2},\qquad |A^{*}|=|A|^{n-1}=2^{2}=4 .
\]

（1）
\[
A^{*}-3A^{-1}=A^{*}-\frac{3}{2}A^{*}=-\frac{1}{2}A^{*},
\]
故
\[
|A^{*}-3A^{-1}|=\left(-\frac{1}{2}\right)^{3}|A^{*}|=-\frac{1}{8}\times4=-\frac{1}{2}.
\]
注意数乘是 $3$ 阶，提出 $-\frac12$ 时是三次方，不是一次方。

（2）由 $(A^{*})^{-1}=A/|A|=A/2$，得
\[
(A^{*})^{-1}+A=\frac{A}{2}+A=\frac{3}{2}A,
\]
故
\[
|(A^{*})^{-1}+A|=\left(\frac{3}{2}\right)^{3}|A|=\frac{27}{8}\times2=\frac{27}{4}.
\]

（3）这是两块都在副对角线上的分块行列式，$m=n=3$，由第 6 题的公式
\[
\begin{vmatrix}0&A\\A^{*}&0\end{vmatrix}=(-1)^{3\times3}|A|\,|A^{*}|=-1\times2\times4=-8 .
\]

\boxed{\textbf{答案：}\text{（1）}-\frac{1}{2},\ \text{（2）}\frac{27}{4},\ \text{（3）}-8}

\else \answerarea[8] \fi%

\textbf{9.} 设 $A$ 是 $n$ 阶实方阵，$n$ 为奇数，且 $A^{T}=-A$（反对称矩阵）。证明 $|A|=0$，并用一个三阶例子验证。

\ifshowsolutions%
\boxed{\textbf{解}}
\textbf{思路：} 把 $|A|$ 用两种方式写成同一个数：一方面转置不改变行列式，另一方面 $A^{T}=-A$ 提出负号时要用 $|-A|=(-1)^{n}|A|$，奇偶性使两式相互矛盾。

由转置不变性与 $A^{T}=-A$：
\[
|A|=|A^{T}|=|-A|=(-1)^{n}|A|.
\]
因为 $n$ 为奇数，$(-1)^{n}=-1$，所以 $|A|=-|A|$，即 $2|A|=0$，从而 $|A|=0$。

\textbf{三阶例子：} 取
\[
A=\begin{pmatrix}0&1&2\\-1&0&3\\-2&-3&0\end{pmatrix},
\]
显然 $A^{T}=-A$。按第一行展开：
\[
|A|=0\cdot(0+9)-1\cdot(0+6)+2\cdot(3-0)=-6+6=0,
\]
与结论一致。

\textbf{评注：} 偶数阶反对称矩阵的行列式一般不为零，例如 $\begin{pmatrix}0&1\\-1&0\end{pmatrix}$ 的行列式为 $1$。所以“奇数阶”这个条件不能去掉。

\boxed{\textbf{答案：}|A|=0}（证明见上）

\else \answerarea[8] \fi%

\subsection{真题风格与综合训练}

\textbf{10.}（选择题）设 $A$ 为 $n$ 阶方阵，满足 $A^{2}=A$ 且 $A\neq I$（$I$ 为单位矩阵），则 $|A|=$\\
(A) $0$ \quad (B) $1$ \quad (C) $-1$ \quad (D) 不能确定

\ifshowsolutions%
\boxed{\textbf{解}}
\textbf{思路：} 对矩阵等式取行列式，得到关于 $|A|$ 的方程；再用“$|A|\neq0$ 等价于 $A$ 可逆”排除其中一个根。

对 $A^{2}=A$ 两边取行列式，利用乘法公式：
\[
|A|^{2}=|A|\quad\Longrightarrow\quad |A|\bigl(|A|-1\bigr)=0,
\]
所以 $|A|=0$ 或 $|A|=1$。

若 $|A|=1\neq0$，则 $A$ 可逆，在 $A^{2}=A$ 两边左乘 $A^{-1}$ 得 $A=I$，与题设 $A\neq I$ 矛盾。故只能 $|A|=0$。

也可举例说明 $|A|=0$ 确实能取到：$A=\begin{pmatrix}1&0\\0&0\end{pmatrix}$ 满足 $A^{2}=A$、$A\neq I$，且 $|A|=0$。

\textbf{评注：} 条件 $A\neq I$ 不可省。若允许 $A=I$，则 $|A|$ 还可能等于 $1$，答案就变成 (D) 了；题设排除了这种情形，所以答案唯一。

\boxed{\textbf{答案：}(A)}

\else \answerarea[4] \fi%

\textbf{11.}（填空题）计算
\[
D=\begin{vmatrix}
1&2&3&4\\
2&3&4&1\\
3&4&1&2\\
4&1&2&3
\end{vmatrix}=\underline{\qquad\qquad}.
\]

\ifshowsolutions%
\boxed{\textbf{解}}
\textbf{思路：} 观察每一行的元素都是 $1,2,3,4$ 的一个轮换，每行之和都是 $10$，每列之和也都是 $10$。把第 $2,3,4$ 列都加到第 $1$ 列，第一列就变成常数 $10$，可提出公因子，再消去第一列的其余元素。

先把 $c_2+c_3+c_4$ 加到 $c_1$：
\[
D=\begin{vmatrix}
10&2&3&4\\
10&3&4&1\\
10&4&1&2\\
10&1&2&3
\end{vmatrix}
=10\begin{vmatrix}
1&2&3&4\\
1&3&4&1\\
1&4&1&2\\
1&1&2&3
\end{vmatrix}.
\]
再作 $r_2-r_1,\ r_3-r_1,\ r_4-r_1$：
\[
D=10\begin{vmatrix}
1&2&3&4\\
0&1&1&-3\\
0&2&-2&-2\\
0&-1&-1&-1
\end{vmatrix}.
\]
按第 $1$ 列展开（只有 $a_{11}=1$）：
\[
D=10\begin{vmatrix}
1&1&-3\\
2&-2&-2\\
-1&-1&-1
\end{vmatrix}
=10\Bigl[1(2-2)-1(-2-2)+(-3)(-2-2)\Bigr]=10\times16=160 .
\]
其中括号里三项分别为 $0$、$4$、$12$，合计 $16$。

\textbf{评注：} 这类“行和相等”的行列式（循环矩阵）都用“把各列加到第一列”起手；若对行操作更顺手，也可把各行加到第一行，效果相同。

\boxed{\textbf{答案：}D=160}

\else \answerarea[8] \fi%

\textbf{12.}（解答题）设 $a$ 为实数，讨论 $a$ 取何值时方程组
\[
\begin{cases}
ax_1+x_2+x_3=1,\\
x_1+ax_2+x_3=1,\\
x_1+x_2+ax_3=1
\end{cases}
\]
有唯一解；在有唯一解时用克莱姆法则求出解；并说明其余情形下方程组的解的情况。

\ifshowsolutions%
\boxed{\textbf{解}}
\textbf{思路：} 三阶系数行列式 $D$ 决定是否有唯一解。观察到每一列的元素之和都等于 $a+2$，把第 $2,3$ 行加到第 $1$ 行即可提出公因子算出 $D$，再讨论 $D=0$ 的两个根。$D\neq0$ 时用克莱姆法则，由对称性只需算一个 $D_j$。

把第 $2,3$ 行都加到第 $1$ 行，第一行变为 $\bigl(a+2,\ a+2,\ a+2\bigr)$，提出公因子 $a+2$：
\[
D=(a+2)\begin{vmatrix}
1&1&1\\
1&a&1\\
1&1&a
\end{vmatrix}.
\]
再作 $r_2-r_1,\ r_3-r_1$：
\[
D=(a+2)\begin{vmatrix}
1&1&1\\
0&a-1&0\\
0&0&a-1
\end{vmatrix}=(a+2)(a-1)^{2}.
\]

\textbf{情形一：$a\neq1$ 且 $a\neq-2$。} 此时 $D=(a+2)(a-1)^{2}\neq0$，方程组有唯一解。由对称性 $D_1=D_2=D_3$，计算 $D_1$（把第 $1$ 列换成常数列 $(1,1,1)^{T}$）：
\[
D_1=\begin{vmatrix}
1&1&1\\
1&a&1\\
1&1&a
\end{vmatrix}
=1(a^{2}-1)-1(a-1)+1(1-a)=(a-1)^{2}.
\]
于是
\[
x_1=x_2=x_3=\frac{D_1}{D}=\frac{(a-1)^{2}}{(a+2)(a-1)^{2}}=\frac{1}{a+2}.
\]
这与“三个方程相加得 $(a+2)(x_1+x_2+x_3)=3$，再由对称性 $x_1=x_2=x_3$”得到的结论一致。

\textbf{情形二：$a=1$。} 三个方程都退化为 $x_1+x_2+x_3=1$，系数矩阵秩为 $1$，有无穷多解：
\[
(x_1,x_2,x_3)=(1-s-t,\ s,\ t),\quad s,t\in\mathbb{R}.
\]

\textbf{情形三：$a=-2$。} 方程变为
\[
-2x_1+x_2+x_3=1,\quad x_1-2x_2+x_3=1,\quad x_1+x_2-2x_3=1 .
\]
三式相加，左边为 $0$，右边为 $3$，矛盾，故无解。

\textbf{评注：} $a=1$ 时 $D=0$ 但方程组有解（不唯一）；$a=-2$ 时 $D=0$ 且无解。这说明“$D=0$”只能说明没有唯一解，具体是无穷多解还是无解，必须回到方程本身判断。

\boxed{\textbf{答案：}a\neq1,-2\text{ 时 }x_1=x_2=x_3=\frac{1}{a+2};\ a=1\text{ 时无穷多解};\ a=-2\text{ 时无解}}

\else \answerarea[10] \fi%
